%---------------------------------------------------------------%
% preamble.tex
%
% Packages and macros for the NE 806 notes, loaded by main.tex.
% Formatting (chapter style, title page) lives in format106.sty.
%---------------------------------------------------------------%

\usepackage{bm,xcolor,cancel}
% This gives a bit more white space...possibly desirable.
%\OnehalfSpace

\DefineNamedColor{named}{mitred}    {rgb}{0.6,0.2,0.2}
\DefineNamedColor{named}{mitgray}   {rgb}{0.4,0.4,0.4}
\DefineNamedColor{named}{darkgray}   {cmyk}{0,0,0,0.90}

\usepackage{format106}
\usepackage{mathtools,array}
\usepackage[shortlabels]{enumitem}
\usepackage{tikz}
\usetikzlibrary{arrows.meta,calc,decorations.pathreplacing,positioning,shapes.geometric}
\usepackage{hyperref} 
\newcommand{\cov}{{$\, \mathbf{C}_{\sigma\sigma} \,$}}
\newcommand{\transp}{\mathsf{T}}


\newcommand{\ISOA}  [2]{$\prescript{#2}{}{\text{#1}}^{}_{}$}
\newcommand{\ISOAZ} [3]{$\prescript{#2}{#3}{\text{#1}}^{}_{}$}
\newcommand{\ISOAZN}[4]{$\prescript{#2}{#3}{\text{#1}}^{}_{#4}$}
\newcommand{\isoA}  [2]{\prescript{#2}{}{\text{#1}}^{}_{}}
\newcommand{\isoAZ} [3]{\prescript{#2}{#3}{\text{#1}}^{}_{}}
\newcommand{\isoAZN}[4]{\prescript{#2}{#3}{\text{#1}}^{}_{#4}}
             
\newcommand{\keff}{$k_{\text{eff}}$ }
\newcommand{\ie}{\textit{i.e. }}
\newcommand{\eg}{\textit{e.g. }}

\newcommand{\LECTURE}[1]{Lecture~\ref{lec:#1}}
%---------------------------------------------------------------%
% Transcription markup
%
% The nuclear data lectures were transcribed from handwritten notes
% and carry markup tying them to the source scans.  The draft build
% (the default) shows it; the student build hides it.  Build the
% student version with
%   latexmk -jobname=main-student -usepretex='\def\studentbuild{}'
%---------------------------------------------------------------%
\newif\ifdraftnotes
\ifdefined\studentbuild \draftnotesfalse \else \draftnotestrue \fi

\ifdraftnotes
\newcommand{\sourcepage}[1]{%
  \par\medskip\noindent
  \textcolor{black!45}{\rule{0.31\linewidth}{0.4pt}\hfill\textsc{Source page #1}\hfill\rule{0.31\linewidth}{0.4pt}}%
  \par\smallskip
}
\newcommand{\marginnoteinline}[1]{\par\smallskip\noindent\textit{Margin note: #1}\par\smallskip}
\newcommand{\editorial}[1]{\textcolor{blue!55!black}{\textit{[Editorial note: #1]}}}
\newcommand{\uncertain}[1]{\textcolor{red!65!black}{\textit{[uncertain: #1]}}}
\else
\newcommand{\sourcepage}[1]{}
\newcommand{\marginnoteinline}[1]{}
\newcommand{\editorial}[1]{}
\newcommand{\uncertain}[1]{}
\fi

\providecommand{\CN}{\mathrm{CN}}
\providecommand{\dd}{\,\mathrm{d}}
\providecommand{\ii}{\mathrm{i}}
\providecommand{\ee}{\mathrm{e}}
\renewcommand{\unit}[1]{\,\mathrm{#1}}
\newcommand{\ketc}[1]{\lvert #1\rangle}
\newcommand{\brac}[1]{\langle #1\rangle}
